\section{Zaključak}

Analiza pokazuje da algoritamska manipulacija pažnjom nije izdvojen mehanizam, već spoj ekonomskog, tehničkog i psihološkog sistema. U ekonomiji pažnje korisnikova pažnja postaje ograničen resurs, a u nadzornom kapitalizmu njegovo ponašanje postaje izvor podataka i predviđanja.

Tehničke metode kao što su collaborative filtering, reinforcement learning i A/B testiranje omogućavaju platformama da iz korisničkog ponašanja uče šta zadržava pažnju i da to znanje stalno primenjuju na novo prilagođavanje sadržaja. Metrike poput klikova, vremena zadržavanja, dubine skrolovanja i retention-a pretvaraju pažnju u merljivu vrednost, a ono što se meri počinje da određuje ono što se optimizuje.

Psihološki mehanizmi pokazuju zašto takva optimizacija deluje. Promenljive nagrade, beskonačni feed, autoplay, notifikacije, socijalna validacija i FOMO ne deluju kao otvorena prinuda, već kao okruženje u kome korisnik sve teže proizvodi prekid. Problem zato ne može biti sveden samo na individualnu slabost ili lošu samokontrolu.

Polazna pretpostavka rada time se potvrđuje: manipulacija pažnjom nastaje upravo u spoju poslovnog interesa, algoritamske optimizacije, merenja ponašanja i psihološkog dizajna digitalnog okruženja. Platforme ne moraju nužno da žele štetu korisniku da bi proizvodile štetne posledice. Dovoljno je da sistem dosledno nagrađuje ono što korisnika duže zadržava, češće vraća i čini predvidljivijim.

Zato pitanje digitalne pažnje nije samo pitanje lične discipline, već pitanje društvene odgovornosti tehnoloških sistema. Ako je pažnja korisnika sirovina, ponašanje proizvod, a zadržavanje cilj, onda se autonomija korisnika ne može štititi samo savetom da „manje koristi telefon”. Potrebno je prepoznati da problem nije samo u korisniku koji ostaje, već i u sistemu koji je napravljen da ga zadrži.